\chapter{Design}
\label{chapter:design}

\section{Design approach - secure by design, SAVOIR-compatible}
\label{section:design-approach}

The main contribution of this work is the proposed Secure Boot SW design, and therefore particular care was taken to consider all relevant information. Following a review of existing standards, research, and recommendations from the space segment, a specification document was developed that could serve as an extension to the existing SAVOIR framework. The objective was to achieve a secure-by-design solution, with security treated as a primary design consideration rather than an afterthought. This section describes how the findings and recommendations presented in \autoref{chapter:bg} were incorporated into the proposed design. 

\subsection{Existing material}
\label{subsection:design-approach-existing}

The specification draws on two preceding chapters. \autoref{chapter:bg} established the technical foundations: how space BSW is structured, what the SAVOIR FCIS framework already mandates, how A/B partitioning and over-the-air updates work, and what the cryptographic literature recommends for long-duration systems including the transition to post-quantum algorithms. \autoref{chapter:threatModel} then defined the attacker scenarios the design must defend against: a physical adversary with pre-launch hardware access, and a remote adversary operating through the communications and update channels after deployment.

Together these inputs shaped three core design decisions. First, SAVOIR already requires the BSW to reside in write-protected Boot Memory and defines integrity checks on ASW images; the gap is that those checks are CRC-based and provide no authenticity guarantee against a deliberate adversary. The specification closes this gap by adding cryptographic signature verification on top of the existing integrity checks, leaving the SAVOIR structure otherwise intact. Second, A/B partitioning is already present in space BSW; the specification extends it with rollback protection to prevent an attacker from exploiting the update channel to reinstall a vulnerable older version. Third, the background survey highlighted that missions with lifetimes of ten years or more face a realistic risk that classical signature algorithms may be weakened or broken before end of missions due to advances in quantum computing. The specification therefore requires that the digital signature algorithm provides resistance against quantum computational attacks.


\section{Secure boot specification (SAVOIR-GS-002S extension)}
\label{section:design-spec}

In order to make the requirements easy to integrate, the same language and structure as existing SAVOIR documents were used. Additionally, standardised keywords are employed to indicate requirement levels, according to RFC 2119 \cite{rfc2119}. Based on the research, all of SAVOIR-GS-002 can be kept as-is, with additional requirements added on top. The one exception is \textbf{SAVOIR.BOOTSW.BEF.22}, which permits the Fast Boot Path to execute a subset of the nominal sequence steps. Because digital signature verification (\textbf{SEC.03}) is a newly introduced step, BEF.22 must be constrained so that signature verification cannot be placed in the skippable subset; this constraint is captured as requirement \textbf{SAVOIR.BOOTSW.BEF.22S} in \autoref{appendix:savoir-secure-boot}, which lists the full formal requirements. This section explains the reasoning behind them in more detail.

\subsection{Establishing a hardware root of trust}

The entire Secure Boot chain depends on a fundamental question: what is the first piece of software that can be trusted, and on what basis? As the Boot SW is the first software executed during system startup, it cannot verify its own integrity through software mechanisms. An attacker capable of modifying the Boot SW could alter its initial instructions and bypass all subsequent verification steps. Consequently, the root of trust cannot reside in software and must instead be established through hardware.

There are two hardware mechanisms that address this. First, the BSW itself must reside in a memory region that is physically immutable at runtime (a write-protected or one-time-programmable flash sector) so that no software, including compromised ASW, can overwrite it. Second, the public key used to verify the ASW images must be stored in an  immutable region as well. Although an adversary who reads the public key gains no useful information, one who can replace it could substitute their own key and sign arbitrary images. Immutable storage disables this attack path entirely. If the immutable storage is limited, a hash of the public key could be stored instead with the public key being included as part of the ASW images.

This approach is consistent with NIST SP~800-193 \cite{NIST-sp-800-193}, which requires that Roots of Trust ``shall either be immutable or protected using mechanisms which ensure all RoTs and CoTs remain in a state of integrity'' (§4.1.1) and that ``the write protection of field non-upgradable memory shall not be modifiable'' (§4.2.2). The standard further specifies that RoT functions ``shall be resistant to any tampering attempted by software running under, or as part of, the operating system on the host processor''~--- a property achieved here by completing all BSW verification before any ASW code is executed.

These two constraints together~--- immutable Boot SW and immutable public key~--- constitute the hardware Root of Trust from which all subsequent verification steps derive their trustworthiness (\textbf{SEC.01-02}). All higher-level security guarantees in this specification are contingent on this foundation holding as well as the private key not becoming compromised.

\subsection{Authenticated execution}

With a trusted starting point established, the specification allows for the trust to be forwarded to the ASW using digital signature verification. Before transferring execution to any ASW image, the Boot SW must verify its digital signature across all image sections including the metadata. A valid signature provides two guarantees simultaneously: integrity and authenticity, ensuring the image has not been modified after signing and that it was signed by a trusted party.

NIST SP~800-193 \cite{NIST-sp-800-193} formalises this as an \emph{Authenticated Update Mechanism} (§4.2.1.1), requiring that ``firmware update images shall be signed using an approved digital signature algorithm'' and that ``the digital signature of a new or recovery firmware update image shall be verified by an RTU or a CTU prior to the non-volatile storage completion of the update process''. The specification also mandates \emph{non-bypassability} (§4.2.1.3): there must be no development or diagnostic interface, architectural feature, or low-level vulnerability that allows an attacker to modify firmware outside the authenticated update mechanism.

The consequence of a failed verification is deliberately strict: the Boot SW must reject the image and prevent execution entirely (\textbf{SEC.03-04}). No fallback exists that permits unsigned or unverified code to run, since any exception would present a path for adversaries to take advantage of. The behaviour on failure is instead handled separately through the rollback and reporting mechanisms described below.

\subsection{Post-quantum cryptographic support}

A bootloader deployed on a long-duration mission may remain in operation well beyond the point at which its original cryptographic algorithms are considered secure. As discussed in \autoref{section:bg-pq}, classical algorithms such as RSA or ECDSA are vulnerable to sufficiently powerful quantum computers, and the transition to post-quantum cryptography is expected to be complete by 2035. For missions with lifetimes of ten years or more, there is therefore a realistic risk that the signature algorithm used to verify the ASW image at boot time may be weakened or broken before end of mission.

The specification addresses this by requiring that the digital signature algorithm used for ASW image verification provides resistance against quantum computational attacks (\textbf{SEC.05}). Algorithm selection remains mission-configurable, and a single algorithm is sufficient provided it offers resistance against both classical and quantum adversaries — hash-based schemes such as LMS/HSS satisfy this property by design.

\subsection{Rollback protection}

A/B partitioning, as described in \autoref{subsection:bg-bootloader-ABandOTA}, is already included in space BSW and provides recovery from failed updates. However, it introduces a risk where an attacker with access to the update mechanism could deliberately install an older, more vulnerable version of the ASW, undoing a security patch. Rollback protection closes this attack path. NIST SP~800-193 \cite{NIST-sp-800-193} explicitly identifies this threat in §4.2.1.3, stating that ``unauthorized rollback could allow an attacker to restore a vulnerable firmware image, which in turn could allow the attacker to damage the device or inject malware'', and requires that authenticated update mechanisms ``be capable of preventing unauthorized updates of the device firmware to an earlier authentic version that has a security weakness''.

Each ASW image carries a strictly positive integer version identifier. The Boot SW maintains a monotonic counter in protected non-volatile memory (which can be updated by the BSW only), recording the highest version that has been executed. Execution is then permitted only for images whose version falls within the inclusive range $[N, N - X]$, where N is the counter value and X is the mission-configured rollback window size (\textbf{SEC.06}).

A window parameter allows for greater flexibility based on the mission-specific needs. A window of zero would provide the strongest security but would eliminate any autonomous recovery from a defective update, requiring ground intervention for every rollback. A window of one permits recovery to the immediately preceding version while still preventing an attacker from reaching arbitrarily old and vulnerable images. A larger window size might also be applicable in specific cases. Storing the counter in protected non-volatile memory ensures it persists across resets and cannot be manipulated by the ASW, which is essential given that a compromised ASW must not be able to lower its own minimum version floor.

\subsection{Failure visibility}

In terrestrial systems, a boot failure is immediately visible. On a spacecraft, the only diagnostic channel is telemetry and the window for ground intervention might be limited or delayed. A silent failure could be indistinguishable from a hardware fault, leading to incorrect corrective action.

For this reason, the design requires the Boot SW to report secure boot failures and their specific causes through the Boot Report or telemetry mechanisms already defined in SAVOIR (\textbf{SEC.07}). Reportable causes include signature verification failure, use of an unrecognised or revoked key and version window violations. This gives ground operators the information needed to distinguish a genuine attack from a corrupted image or a configuration error and to take corrective action such as uploading a replacement image. This requirement is aligned with NIST SP~800-193 \cite{NIST-sp-800-193} §3.6.5, which prescribes event logging for firmware security events for forensic analysis and audit purposes, and with §4.3.1, which requires that detection mechanisms ``should be capable of creating notifications of corruption'' and ``should be capable of logging events when corruption is detected''.


\section{Key design trade-offs and discarded options}
\label{section:design-tradeoffs}

Not every security/safety mechanism considered during the design process made it into the final specification. Some were discarded because they conflicted with other requirements, others because their added complexity was not justified by the security benefit they provided.

One notable candidate was the concept of a golden image \cite{cubesat_flight_software}~--- a read-only, pre-launch copy of the ASW stored in a dedicated partition and used as a last resort if all other partitions are corrupted or unbootable. This idea improves the systems safety as it provides a recovery path that does not depend on the update mechanism functioning correctly. However, it goes directly against rollback protection. A golden image is by definition fixed at the version present before launch, meaning that as the mission progresses and newer ASW versions are deployed, the golden image will inevitably fall outside the permitted rollback window. Exempting it from the version check would require a special case path in the Boot SW, which introduces both complexity and a potential bypass vector, allowing an attacker who can trigger the golden image path to effectively force execution of an old, potentially more vulnerable version of the ASW. Making the golden image sector read-only would prevent the adversary from replacing its contents, but would not prevent them from exploiting the existing vulnerabilities in that version. Given that the rollback window provides a controlled recovery mechanism and the ground operators have the ability to upload new images via the Standby mode, the golden image concept was discarded in favour of simpler and more consistent design.

\section{How does this address the threat model?}
\label{section:design-threatModel}

The two attacker scenarios defined in \autoref{chapter:threatModel} place different demands on the design and it is worth examining how the specification responds to each.

For a physical attacker, the critical window is the period between manufacturing and launch, during which the hardware may be accessible and debug interfaces may still be enabled. The immutable storage requirements (\textbf{SEC.01-02}) are the primary defences here. Since the BSW and its public key(s) reside in write-protected memory, a physical attacker cannot modify the verification logic or substitute a different key without triggering a full device erase, which would be immediately apparent. In practice, the correct pre-launch procedure is to reset the ASW partitions and upload the mission images immediately before the system is finalised, at which point debug access should be revoked. Since the BSW and public key(s) are immutable by that point, no modification made during earlier integration phases can survive into flight.

For a remote attacker operating after launch, the threat surface shifts to the update and communication channels. Here the signature verification requirements (\textbf{SEC.03-04}) form the central control: regardless of how an attacker delivers a modified firmware image (whether by exploiting a software vulnerability, intercepting a telecommand or abusing the update mechanism), the Boot SW will refuse to execute it unless it carries a valid signature from the holder of the private key. Since the private key never leaves the ground segment, this guarantee holds as long as the key itself is not compromised. The rollback protection (\textbf{SEC.06}) adds a further constraint, ensuring that even a validly signed but deliberately outdated image cannot be used to reintroduce previously patched vulnerabilities. Together, these mechanisms mean that an attacker without access to the private key has no viable path to persistent code execution on the system, regardless of what access they have achieved at the application level.